\section{Code Snippets}

The following sections contain key code excerpts from the RayCastED model
implementation, organized by architectural component.

\subsection{Architecture Configuration}

The following YAML defines the RayCastED model backbone, neck, and detection
head.

\begin{lstlisting}[language=XML]
nc: 5
end2end: True
reg_max: 1
depth_multiple: 1.0
width_multiple: 1.0
scales:
  n: [0.50, 0.25, 1024]
  s: [0.50, 0.50, 1024]
  m: [0.50, 1.00, 512]
  l: [1.00, 1.00, 512]
  x: [1.00, 1.50, 512]

backbone:
  - [-1, 1, ResoConv, [64]]
  - [-1, 1, ResoConv, [128]]
  - [-1, 2, C3k2, [256, False, 0.25]]
  - [-1, 1, ResoConv, [256]]
  - [-1, 2, C3k2, [512, False, 0.25]]
  - [-1, 1, ResoConv, [512]]
  - [-1, 2, C3k2, [512, True]]
  - [-1, 1, SPPF, [512, 5, 3, True]]
  - [-1, 2, C2PSA, [512]]

head:
  - [8, 1, nn.Upsample, [None, 2, "nearest"]]
  - [[-1, 4], 1, Concat, [1]]
  - [-1, 2, C3k2, [256, True]]

  - [-1, 1, nn.Upsample, [None, 2, "nearest"]]
  - [[-1, 2], 1, Concat, [1]]
  - [-1, 2, C3k2, [128, True]]

  - [-1, 1, ResoConv, [128]]
  - [[-1, 11], 1, Concat, [1]]
  - [-1, 2, C3k2, [256, True]]

  - [-1, 1, ResoConv, [256]]
  - [[-1, 8], 1, Concat, [1]]
  - [-1, 2, C3k2, [512, True]]

  - [[14, 17, 20], 1, RayCastDetect, [nc, reg_max, end2end]]
\end{lstlisting}

\subsection{Detection Head Architecture}

\subsubsection{Dual-Classification Head}

\begin{lstlisting}[language=Python]
# head.py:189-219 -- RayCastDetect.__init__ excerpt

# Separate o2o box and cls heads via deepcopy.
# Shared cls head caused 11.5x overprediction: o2m's dense positives
# taught it to fire high scores for many anchors per GT.
self.one2one_cv2 = copy.deepcopy(self.cv2)
self.one2one_cv3 = copy.deepcopy(self.cv3)

if hierarchical_cls:
    # Hierarchical: binary (fg/bg) + class (cell type) heads.
    # Binary head: ALL anchors get gradient -> strong fg/bg signal.
    # Class head: only fg anchors -> simpler inter-class task.
    self.one2one_cv3_binary = copy.deepcopy(self.cv3)
    for seq in self.one2one_cv3_binary:
        seq[-1] = nn.Conv2d(seq[-1].in_channels, 1, 1)
    self.one2one_cv3_class = copy.deepcopy(self.cv3)
\end{lstlisting}

\subsubsection{RayCastDetect Forward Head}

\begin{lstlisting}[language=Python]
# head.py:221-258 -- RayCastDetect.forward_head

poly = torch.cat([box_head[i](x[i]).view(bs, self.raycast_dim, -1)
                  for i in range(self.nl)], dim=-1)

if cls_head_binary is not None and cls_head_class is not None:
    binary_in = [xi.detach() for xi in x] if self.hierarchical_cls_detach else x
    binary_scores = torch.cat(
        [cls_head_binary[i](binary_in[i]).view(bs, 1, -1)
         for i in range(self.nl)], dim=-1)
    class_scores = torch.cat(
        [cls_head_class[i](x[i]).view(bs, self.nc, -1)
         for i in range(self.nl)], dim=-1)
    result = dict(boxes=poly, binary_scores=binary_scores,
                  class_scores=class_scores, feats=x)
else:
    scores = torch.cat([cls_head[i](x[i]).view(bs, self.nc, -1)
                        for i in range(self.nl)], dim=-1)
    result = dict(boxes=poly, scores=scores, feats=x)
\end{lstlisting}

\begin{lstlisting}[language=Python]
# head.py:269-294 -- RayCastDetect.forward

preds = self.forward_head(x, **self.one2many)
if self.end2end:
    x_detach = [xi.detach() for xi in x]
    one2one = self.forward_head(x_detach, **self.one2one)
    preds = {'one2many': preds, 'one2one': one2one}

if self.training and self.end2end and self.inter_scale_competition and self.nl > 1:
    one2one = preds['one2one']
    if self.hierarchical_cls and 'binary_scores' in one2one:
        binary_prob = one2one['binary_scores'].sigmoid()
        adjusted = self._apply_inter_scale_competition(
            binary_prob, one2one['feats'])
        adjusted = adjusted.clamp(min=1e-7, max=1 - 1e-7)
        one2one['binary_scores'] = torch.log(adjusted / (1 - adjusted))
    else:
        scores_prob = one2one['scores'].sigmoid()
        adjusted = self._apply_inter_scale_competition(
            scores_prob, one2one['feats'])
        adjusted = adjusted.clamp(min=1e-7, max=1 - 1e-7)
        one2one['scores'] = torch.log(adjusted / (1 - adjusted))
    return preds
\end{lstlisting}

\subsubsection{Inter-Scale Competition}

\begin{lstlisting}[language=Python]
# head.py:296-356 -- RayCastDetect._apply_inter_scale_competition

bs, nc, N_total = scores.shape
spatial_shapes = [(f.shape[2], f.shape[3]) for f in feats]
anchor_counts = [h * w for h, w in spatial_shapes]
H_p2, W_p2 = spatial_shapes[0]

scale_scores = scores.split(anchor_counts, dim=-1)

scale_spatial = []
for ss, (h, w) in zip(scale_scores, spatial_shapes):
    scale_spatial.append(ss.view(bs, nc, h, w))

# Upsample all scales to P2 resolution
upsampled = []
for si, ss in enumerate(scale_spatial):
    if si == 0:
        upsampled.append(ss)
    else:
        upsampled.append(F.interpolate(
            ss, size=(H_p2, W_p2), mode='bilinear', align_corners=False))

stacked = torch.stack(upsampled, dim=2)

if self.inter_scale_temperature != 1.0:
    stacked = stacked / self.inter_scale_temperature

# Softmax competition across the scale dimension
competition_weights = F.softmax(stacked, dim=2)

# Apply per-scale suppression weights
result_parts = []
for si in range(len(scale_spatial)):
    cw = competition_weights[:, :, si, :, :]
    if si == 0:
        adjusted = scale_spatial[si] * cw
    else:
        cw_down = F.interpolate(
            cw, size=spatial_shapes[si],
            mode='bilinear', align_corners=False)
        adjusted = scale_spatial[si] * cw_down
    result_parts.append(adjusted.view(bs, nc, -1))

return torch.cat(result_parts, dim=-1)
\end{lstlisting}

\subsection{Task-Aligned Assignment}

\subsubsection{Alignment Metric}

\begin{lstlisting}[language=Python]
# tal.py:159-230 -- RayCastAssigner.get_box_metrics

na = pd_bboxes.shape[-2]
mask_gt_bool = mask_gt.bool().expand(-1, -1, na)
overlaps = torch.zeros([self.bs, self.n_max_boxes, na],
                        dtype=torch.float32, device=pd_bboxes.device)
bbox_scores = torch.zeros([self.bs, self.n_max_boxes, na],
                           dtype=pd_scores.dtype, device=pd_scores.device)

# Gather cls scores per GT class
ind = torch.zeros([2, self.bs, self.n_max_boxes], dtype=torch.long)
ind[0] = torch.arange(end=self.bs).view(-1, 1).expand(-1, self.n_max_boxes)
ind[1] = gt_labels.squeeze(-1)
bbox_scores[mask_gt_bool] = pd_scores[ind[0], :, ind[1]][mask_gt_bool]

if gauss_data is not None:
    topk_idx, _gauss = gauss_data
    k_prefilt = topk_idx.shape[-1]

    if self.use_cls_only:
        overlaps.scatter_(2, topk_idx,
            torch.ones(self.bs, self.n_max_boxes, k_prefilt,
                        dtype=overlaps.dtype, device=overlaps.device))
    else:
        n_rays = pd_bboxes.shape[-1] - 2
        bs_idx = torch.arange(self.bs, device=topk_idx.device)
        bs_idx = bs_idx.view(-1, 1, 1).expand(-1, self.n_max_boxes, k_prefilt)
        pd_block = pd_bboxes[bs_idx, topk_idx, 2:]
        gt_rays = gt_bboxes[:, :, 2:].unsqueeze(2).expand(-1, -1, k_prefilt, -1)

        iou_block = polar_iou_torch(pd_block.reshape(-1, n_rays),
                                     gt_rays.reshape(-1, n_rays))
        iou_block = iou_block.reshape(self.bs, self.n_max_boxes, k_prefilt)
        overlaps.scatter_(2, topk_idx, iou_block.to(overlaps.dtype))

# align = cls^alpha * PolarIoU^beta
align_metric = bbox_scores.pow(self.alpha) * overlaps.pow(self.beta)
return align_metric, overlaps
\end{lstlisting}

\subsubsection{Gaussian Spatial Decay and Dynamic Top-K Selection}

\begin{lstlisting}[language=Python]
# tal.py:236-274 -- RayCastAssigner.get_pos_mask

gauss_data = (self._prefilter_by_distance(gt_bboxes, anc_points, mask_gt)
              if self.radius_scale > 0 else None)

align_metric, overlaps = self.get_box_metrics(
    pd_scores, pd_bboxes, gt_labels, gt_bboxes, mask_gt, gauss_data=gauss_data)

# Apply Gaussian spatial decay: exp(-d^2 / 2*sigma^2)
# Distant anchors get ~0 naturally -- no hard containment filtering needed.
if gauss_data is not None:
    topk_idx, gauss = gauss_data
    current = align_metric.gather(2, topk_idx)
    updated = current * gauss
    align_metric.scatter_(2, topk_idx, updated)

# Dynamic per-GT top-k selection
mask_topk = self.select_topk_candidates(
    align_metric, topk_mask=mask_gt.expand(-1, -1, self.topk).bool())

mask_pos = mask_topk * mask_gt

# STAL: enforce minimum positive assignments per GT
if self.stal_min_positives > 0:
    pos_per_gt = mask_pos.sum(dim=-1)
    valid_gt = mask_gt.any(dim=-1)
    missing = (pos_per_gt < self.stal_min_positives) & valid_gt
    if missing.any():
        distances = torch.cdist(
            gt_bboxes[:, :, :2].float(), anc_points.float())
        _, nearest_idx = distances.min(dim=-1)
        for b in range(self.bs):
            missing_gts = missing[b].nonzero(as_tuple=True)[0]
            for g in missing_gts:
                mask_pos[b, g, nearest_idx[b, g]] = 1

return mask_pos, align_metric, overlaps
\end{lstlisting}

\subsection{Ground-Truth Ray Computation}

\subsubsection{Cramer's Rule Solver for Ray--Polygon Intersection}

\begin{lstlisting}[language=Python]
# raycast_gpu.py:200-243 -- _solve_ray_intersections

cos_d, sin_d = directions
n_polys, n_verts, _ = v_padded.shape
poly_idx = torch.arange(n_polys, device=centroids.device)

# Connect last real vertex to first (close polygon loop)
n_real = mask.sum(dim=1).long()
vjn = torch.empty_like(v_padded)
vjn[:, :-1, :] = v_padded[:, 1:, :]
vjn[poly_idx, n_real - 1] = v_padded[poly_idx, 0]

last_edge = (torch.arange(n_verts, device=centroids.device).unsqueeze(0)
             == (n_real - 1).unsqueeze(1))
edge_mask = mask & (torch.arange(n_verts, device=centroids.device)
                     .unsqueeze(0) < n_real.unsqueeze(1))
edge_mask = edge_mask | last_edge

v1x, v1y = v_padded[:, :, 0].unsqueeze(1), v_padded[:, :, 1].unsqueeze(1)
v2x, v2y = vjn[:, :, 0].unsqueeze(1), vjn[:, :, 1].unsqueeze(1)
cx = centroids[:, 0].unsqueeze(1).unsqueeze(2)
cy = centroids[:, 1].unsqueeze(1).unsqueeze(2)

dx, dy = v2x - v1x, v2y - v1y
cos_exp, sin_exp = cos_d.view(1, -1, 1), sin_d.view(1, -1, 1)

det = -cos_exp * dy + sin_exp * dx
det_valid = det.abs() > 1e-10

# Cramer's rule for ray--edge intersection
t_num = (v1x - cx) * (-dy) + (v1y - cy) * dx
u_num = (v1x - cx) * (-sin_exp) + (v1y - cy) * cos_exp

safe_det = torch.where(det_valid, det, torch.ones_like(det))
t = t_num / safe_det
u = u_num / safe_det

valid = det_valid & (t > 1e-6) & (u >= 0) & (u <= 1)
valid = valid & edge_mask.unsqueeze(1)

dists = torch.where(valid, t, torch.full_like(t, 1e10))
min_dists, _ = dists.min(dim=-1)

return min_dists
\end{lstlisting}

\subsection{Geometric Loss Formulation}

\subsubsection{Rays Loss and Analytical GT Distances}

\begin{lstlisting}[language=Python]
# star_distances_analytical.py:152-217 -- _analytical_ray_distances_paired

def _analytical_ray_distances_paired(centroids_px, vertices, ray_cos, ray_sin):
    device = centroids_px.device
    dtype = centroids_px.dtype
    n_rays = ray_cos.shape[0]
    n_pairs, n_verts, _ = vertices.shape

    ray_cos = ray_cos.to(device=device, dtype=dtype).view(1, 1, n_rays)
    ray_sin = ray_sin.to(device=device, dtype=dtype).view(1, 1, n_rays)

    v1 = vertices
    v2_shift = vertices.roll(-1, dims=1)

    cx = centroids_px[:, 0:1, None]
    cy = centroids_px[:, 1:2, None]

    v1x = v1[:, :, 0:1]
    v1y = v1[:, :, 1:2]
    v2x = v2_shift[:, :, 0:1]
    v2y = v2_shift[:, :, 1:2]

    dx = v2x - v1x
    dy = v2y - v1y

    cos_e = ray_cos.expand(n_pairs, n_verts, n_rays)
    sin_e = ray_sin.expand(n_pairs, n_verts, n_rays)

    det = -cos_e * dy + sin_e * dx
    det_valid = det.abs() > 1e-10

    t_num = (v1x - cx) * (-dy) + (v1y - cy) * dx
    u_num = (v1x - cx) * (-sin_e) + (v1y - cy) * cos_e

    safe_det = torch.where(det_valid, det, torch.ones_like(det))
    t = t_num / safe_det
    u = u_num / safe_det

    valid = det_valid & (t > 1e-6) & (u >= 0) & (u <= 1)

    large_val = torch.full_like(t, 1e10)
    dists = torch.where(valid, t, large_val)
    min_dists, _ = dists.min(dim=1)

    diagonal_px = torch.as_tensor(
        (vertices.max() - vertices.min()) * 2.0,
        device=device, dtype=dtype).clamp(min=256.0)
    min_dists = min_dists.clamp(max=diagonal_px)
    return min_dists
\end{lstlisting}

\begin{lstlisting}[language=Python]
# star_distances_analytical.py:220-244 -- analytical_gt_rays

def analytical_gt_rays(pred_xy_px, gt_vertices, ray_cos, ray_sin, crop_size):
    distances = _analytical_ray_distances_paired(
        pred_xy_px, gt_vertices, ray_cos, ray_sin)
    return distances / crop_size
\end{lstlisting}

\subsubsection{Polar IoU Loss}

\begin{lstlisting}[language=Python]
# loss.py:402-412 -- L_PolarIoU in get_assigned_targets_and_loss

# -log(PolarIoU), clamped to bound gradient magnitude
fg_piou = polar_iou_torch(fg_pred_rays, fg_target_rays_static)
piou_loss = -torch.log(fg_piou.clamp(min=1e-4))
if fg_plb is not None:
    piou_loss = piou_loss * fg_plb
loss[3] = piou_loss.sum() / n_fg
\end{lstlisting}

\subsubsection{1D Discrete Laplacian Smooth Loss}

\begin{lstlisting}[language=Python]
# ops/loss.py:29-56 -- curvature_smoothness_loss_torch

def curvature_smoothness_loss_torch(rays):
    """Circular 2nd-order difference (curvature) penalty.

    L_curv = mean|d_{i+1} - 2*d_i + d_{i-1}|

    A perfect circle has zero penalty. Smooth irregular shapes
    (ellipses, organic nuclei) have low penalty. Only sharp kinks
    and zigzags are penalised -- unlike 1st-order smoothness, which
    also fights legitimate non-round GT shapes.

    Circular indexing wraps at ends so every ray contributes.
    """
    rays_next = torch.roll(rays, shifts=-1, dims=-1)
    rays_prev = torch.roll(rays, shifts=1, dims=-1)
    curvature = rays_next - 2.0 * rays + rays_prev
    return torch.mean(torch.abs(curvature), dim=-1)
\end{lstlisting}

\subsubsection{Training-Time Integration of Geometric Terms}

\begin{lstlisting}[language=Python]
# loss.py:370-423 -- Training-time analytical ray computation and
# final three geometric loss terms

ray_cos, ray_sin = build_ray_directions(
    self.n_rays, device=self.device, dtype=torch.float32)
gt_vertices = _normed_rays_to_vertices(
    gt_centroids_px, gt_rays_norm, crop_size_val, ray_cos, ray_sin)
gt_vertices = gt_vertices.reshape(*gt_bboxes.shape[:2], self.n_rays, 2)

fg_indices = fg_mask.float().nonzero(as_tuple=False)
fg_batch_idx = fg_indices[:, 0]
fg_gt_idx = target_gt_idx[fg_indices[:, 0], fg_indices[:, 1]]
matched_gt_vertices = gt_vertices[fg_batch_idx, fg_gt_idx]

# Recompute GT rays from predicted centroid to GT polygon boundary
analytical_target_rays = analytical_gt_rays(
    fg_pred_xy_px.detach(), matched_gt_vertices,
    ray_cos, ray_sin, crop_size_val)
fg_target_rays = analytical_target_rays.to(dtype=fg_target_rays.dtype)

# L_L1: Uniform MAE on ray distances
if self.log_ray_loss:
    loss_l1 = _log_space_ray_loss(fg_pred_rays, fg_target_rays)
else:
    loss_l1 = (fg_pred_rays - fg_target_rays).abs().mean(-1)
loss[2] = loss_l1.sum() / n_fg

# L_PolarIoU: -log(PolarIoU) -- uses static GT for polygon overlap
fg_piou = polar_iou_torch(fg_pred_rays, fg_target_rays_static)
loss[3] = (-torch.log(fg_piou.clamp(min=1e-4))).sum() / n_fg

# L_smooth: Circular 2nd-order curvature regularisation
smooth_loss = curvature_smoothness_loss_torch(fg_pred_rays)
loss[4] = smooth_loss.sum() / n_fg
\end{lstlisting}
